\documentclass[11pt,a4paper]{article}
\usepackage[margin=2cm,headheight=14pt]{geometry}
\usepackage[T1]{fontenc}
\usepackage[utf8]{inputenc}
\usepackage{lmodern,microtype,booktabs,tabularx,enumitem,fancyhdr}
\usepackage{xcolor}
\definecolor{linkblue}{HTML}{174A80}
\usepackage[colorlinks=true,urlcolor=linkblue,linkcolor=linkblue]{hyperref}
\hypersetup{pdftitle={Research Methods 2026 -- Coding Assignment 2},pdfauthor={Swapnil Singh}}
\pagestyle{fancy}
\fancyhf{}
\fancyhead[L]{Research Methods 2026}
\fancyhead[R]{Coding Assignment 2}
\fancyfoot[C]{\thepage}
\setlength{\parindent}{0pt}
\setlength{\parskip}{5pt}
\setlist[enumerate]{label=(\alph*),leftmargin=*,topsep=2pt,itemsep=3pt,parsep=0pt}
\setlist[itemize]{leftmargin=*,topsep=2pt,itemsep=1pt,parsep=0pt}
\newcommand{\py}[1]{\texttt{\detokenize{#1}}}
\newcommand{\task}[2]{\subsection*{#1\hfill{\normalsize[#2 marks]}}}
\newcommand{\dataurl}{https://raw.githubusercontent.com/machow/nycflights13-py/1d40369400d77379a78641afd119fc302a5a9659/nycflights13/data}

\begin{document}
\begin{center}
{\Large\textbf{Coding Assignment 2}}\\[3pt]
{\large What can a traveller learn from a year of flights?}
\end{center}

A travel editor wants to know how often flights from New York arrived late in 2013.
Give the editor a table they can check. Then pitch a question of your own and
answer it with one figure the editor could publish.

\textbf{Individual take-home assignment}\quad 100 marks $\cdot$ 10\% of the course grade
$\cdot$ about three hours. Use pandas and plotnine from the Day~2 environment.

Open your pull request by \textbf{Monday, 16 November 2026 at 23:59 Vilnius time}.

\subsection*{Get the data}
Download both files into a folder called \py{data} beside your notebook.
Keep the original files so that someone can run your analysis again.

\begin{itemize}
\item \href{\dataurl/flights.csv.zip}{\textbf{Download flights.csv.zip}} for
336,776 scheduled flight records from Newark (EWR), Kennedy (JFK)
and LaGuardia (LGA) in 2013. Keep this file zipped.
\item \href{\dataurl/airlines.csv}{\textbf{Download airlines.csv}} to get the
full airline \py{name} for each \py{carrier} code.
If the CSV opens as text, save it as \py{airlines.csv}.
\end{itemize}

\begin{quote}\ttfamily\small
import pandas as pd\\
flights = pd.read\_csv("data/flights.csv.zip")\\
airlines = pd.read\_csv("data/airlines.csv")
\end{quote}

Departure and arrival delays are measured in minutes against the schedule.
Negative values mean early and zero means on time. The \py{hour} column is the
scheduled departure hour. The data come from the U.S. Bureau of Transportation
Statistics via Hadley Wickham's \emph{nycflights13} (CC0). Michael Chow provides the CSV files.
See the \href{https://nycflights13.tidyverse.org/reference/flights.html}{flight variable guide}
and \href{https://nycflights13.tidyverse.org/reference/airlines.html}{airline lookup guide}.

\task{1. Build the analysis sample}{20}
\begin{enumerate}
\item State what one row represents in each file. Show both shapes, check the
type of \py{arr_delay}, and count exact duplicate flight rows across all original
columns. Explain whether any rows need removing. An airline flight number can
recur on different dates. \hfill [8]
\item Show how many flight records have a known arrival delay and how many have
a missing one. Check that the two counts add up to your flight total.
Create a late-arrival indicator for delays of \textbf{15 minutes or more}.
Keep the outcome missing when the delay is missing. Show how your rule handles
an early arrival, 14 minutes, exactly 15 minutes and a missing value. \hfill [12]
\end{enumerate}

\task{2. Attach the airline names}{10}
Join the lookup to the flight records using \py{carrier}. Check that each lookup
code occurs once, validate the many-to-one join in pandas, and show the row
counts before and after. Report how many flights have an unmatched code.
Use the airline names in your table on the next page.

\newpage
\task{3. Give the editor a table they can check}{30}
\begin{enumerate}
\item Identify the five busiest airlines by their number of scheduled flight
records. Present one row per airline, busiest first, then a final row for
\textbf{all airlines combined}, including those outside the top five.
Show the airline name, scheduled flights, flights with a known arrival delay,
late arrivals, and the percentage late among flights with a known arrival delay.
Give columns reader-friendly headings and round only the displayed percentages
to one decimal place. \hfill [18]
\item Pick one of the five airlines. Filter its underlying flight records and
recompute its numerator and denominator independently of the summary table.
Show the arithmetic for that airline's percentage and check that it agrees
with the table. \hfill [6]
\item Suppose the editor divides the total number of known late arrivals by
\emph{all scheduled flights}. Calculate this alternative for the all-airlines
row. Report how much it changes the headline rate in percentage points and
explain what the alternative denominator counts. \hfill [6]
\end{enumerate}

\task{4. Pitch and illustrate your own finding}{40}
Choose a reader and a question you can answer from these records. Go beyond
redrawing the five-airline table. For example, how does the chance of a late
arrival change from morning to evening at each airport? How often does a flight that leaves at least
15 minutes late still arrive by its scheduled time? You may pursue a different question.

\begin{enumerate}
\item \textbf{The pitch.} Name your reader and the comparison they care about.
Define the measure, the flights you include and any cutoffs you choose. Give
the number of flights behind the figure. \hfill [8]
\item \textbf{The figure.} Make one figure in plotnine. Choose the geometry,
variables, grouping and scale that bring out your comparison. Give it a headline
supported by the data, label the units and make the groups easy to distinguish
at the size a reader will see. \hfill [14]
\item \textbf{The edit.} Keep a first version in the notebook. Change one design
decision that affects how a reader sees the evidence. Show the final version
and explain in two sentences what is easier to see and what detail, if any,
is lost. A simple figure can earn full marks. \hfill [8]
\item \textbf{The caption.} In 60--100 words, answer your question with a numerical
comparison calculated in the notebook. Name the source, period and one specific
limit on what the comparison can establish. \hfill [10]
\end{enumerate}

\subsection*{What to submit}
Open a pull request to \py{main} from your own branch. In
\py{submissions/exercise-02/YOUR-USERNAME/}, include \py{exercise_02.ipynb}
with the table, figure drafts and caption, the final \py{figure.png}, and a
short \py{README.md} with the download and run instructions. Keep the two raw
downloads on your computer. Restart the kernel and run all cells in order to
check that the notebook recreates the table and saves the final figure.
\end{document}
